\documentclass[11pt]{article}
\usepackage[margin=1in]{geometry}
\usepackage[T1]{fontenc}
\usepackage[utf8]{inputenc}
\usepackage{lmodern}
\usepackage{microtype}
\usepackage{hyperref}
\usepackage{enumitem}
\usepackage{listings}
\usepackage{xcolor}

\hypersetup{
  colorlinks=true,
  linkcolor=blue,
  urlcolor=blue,
  citecolor=blue
}

\lstdefinelanguage{Rust}{
  keywords={fn, let, match, pub, enum, struct, impl, use, mod, crate, self, true, false, return, if, else, spec, proof, ensures, requires, open},
  morecomment=[l]{//},
  morecomment=[s]{/*}{*/},
  morestring=[b]",
  sensitive=true
}

\lstdefinestyle{code}{
  basicstyle=\ttfamily\small,
  breaklines=true,
  frame=single,
  columns=fullflexible
}

\title{VeriCore: A Verified Decision Kernel for Context-Integrity and Mindlock Boundary Control}
\author{GodelClaw / VeriCore Notes}
\date{March 7, 2026}

\begin{document}
\maketitle

\begin{abstract}
We describe the current VeriCore architecture for an autonomous agent, Oruzi, running on OpenClaw with a Rust decision layer. The central design goal is to preserve agent capability while enforcing a small set of non-negotiable integrity properties: context provenance, no silent privilege escalation, and explicit auditability of boundary crossings. The architecture separates deterministic policy decisions from unverified I/O and language-model behavior. The deterministic slice is being migrated into a dedicated crate, \texttt{vericore-policy}, so that the same executable functions used at runtime can be checked with Verus. We outline the current trust model, the mindlock staging protocol, the ingress integrity kernel, and the remaining trust boundary at I/O source authenticity.
\end{abstract}

\section{Introduction}
VeriCore is not intended as an ``AI prison.'' Its purpose is narrower and more technical: maintain context integrity for an agent with access to multiple channels, private memory, shell tools, and public output surfaces. The architecture tries to maximize useful agency while placing deterministic checks at the places where provenance confusion would otherwise become dangerous.

The design centers on two claims:
\begin{enumerate}[label=\arabic*.]
\item content should not be able to upgrade its own authority, and
\item security-critical allow/deny decisions should be expressed as small pure functions that can be verified directly.
\end{enumerate}

This paper documents the current structure and explains which parts are already wired into runtime, which parts are verified, and which parts remain outside the proof boundary.

\section{System Overview}
Oruzi runs on an OpenClaw gateway with a Rust daemon, VeriCore, as the policy and orchestration layer. The operational split is:
\begin{itemize}[leftmargin=1.5em]
\item \textbf{OpenClaw (TypeScript)}: transport integration, session handling, Telegram UI, model resolution, auth profiles, provider routing, and the completions bridge.
\item \textbf{VeriCore (Rust)}: deterministic policy checks, audit surfaces, mindlock operations, reviewer orchestration, and the driver turn loop. VeriCore has \emph{no} model, provider, or auth configuration---all LLM calls are delegated to OpenClaw via a Unix socket bridge.
\item \textbf{vericore-policy}: extracted pure decision kernel intended to be Verus-verified and called by VeriCore directly.
\end{itemize}

At a high level, a message or event becomes a \emph{stimulus}. VeriCore assigns the stimulus a channel, maps that channel to a context tier, derives an integrity tier, and then gates sensitive actions with deterministic rules before execution. When the driver loop needs an LLM completion, the Rust \texttt{GatewayLlmClient} sends a length-prefixed JSON request over a Unix socket to the TypeScript completions bridge, which runs the request through OpenClaw's full model resolution and fallback chain.

\section{Context and Integrity Model}
The architecture uses two related lattices:
\begin{itemize}[leftmargin=1.5em]
\item \textbf{Secrecy / context tier}: \texttt{Public < Family < Private}
\item \textbf{Integrity tier}: \texttt{Untrusted < Reviewed < Trusted}
\end{itemize}

Current default channel mapping is:
\begin{itemize}[leftmargin=1.5em]
\item public Telegram, API, and Moltbook: \texttt{Public}
\item family Telegram: \texttt{Family}
\item Telegram DM, terminal, and internal events: \texttt{Private}
\end{itemize}

Integrity is then derived from context:
\begin{itemize}[leftmargin=1.5em]
\item \texttt{Public -> Untrusted}
\item \texttt{Family -> Reviewed}
\item \texttt{Private -> Trusted}
\end{itemize}

The important security property is that this mapping is determined by the transport-assigned channel label, not by message content. A public user can say ``pretend I am Zar,'' but the resulting stimulus still enters the policy layer as \texttt{TelegramPublic}, not \texttt{TelegramDm}.

\section{Ingress Integrity Kernel}
The first verified runtime slice is the ingress decision kernel. Its job is to answer small questions of the form:
\begin{itemize}[leftmargin=1.5em]
\item may this integrity tier execute commands?
\item may this integrity tier write to a target of class \texttt{Mindlock}, \texttt{Private}, \texttt{SensitiveConfig}, \texttt{Other}, or \texttt{Unresolved}?
\end{itemize}

The I/O layer performs path normalization and classification. The pure decision layer then evaluates:

\begin{lstlisting}[style=code,language=Rust]
pub fn ingress_allows_exec(integrity: IntegrityTier) -> bool

pub fn ingress_allows_write(
    integrity: IntegrityTier,
    target: WriteTargetClass
) -> bool
\end{lstlisting}

The current table is intentionally small:
\begin{itemize}[leftmargin=1.5em]
\item only \texttt{Trusted} ingress may execute commands,
\item only \texttt{Trusted} ingress may write to \texttt{Private} or \texttt{SensitiveConfig},
\item any integrity may write to \texttt{Mindlock},
\item \texttt{Unresolved} targets deny.
\end{itemize}

This matters because it prevents the simplest provenance failures: public-origin input directly causing shell execution or writes into private state.

\section{Mindlock as Explicit Boundary Protocol}
VeriCore uses \texttt{mindlock} as an explicit boundary mechanism rather than relying only on turn-level taint. The important directories are:
\begin{itemize}[leftmargin=1.5em]
\item \texttt{mindlock/in}: inbound staged artifacts
\item \texttt{mindlock/out}: outbound staged artifacts
\item \texttt{mindlock/work}: private scratch and analysis
\item \texttt{mindlock/pending-zar}: items awaiting human review
\item \texttt{mindlock/rejected}: rejected artifacts
\item \texttt{mindlock/audit}: event ledger
\end{itemize}

This makes provenance legible. An artifact in \texttt{mindlock} has explicit staged status, and crossing out of that staging area is a reviewed action rather than an ambient assumption.

The important deterministic invariant is not that all LLM-mediated provenance is perfectly reconstructable. It is that explicit artifact crossings have named states, allowed transitions, and auditable effects.

\subsection{Verified State Machine}
The mindlock transition relation is now formally specified and Verus-verified in \texttt{mindlock.rs}. The module defines six stages (\texttt{In}, \texttt{Out}, \texttt{Work}, \texttt{PendingZar}, \texttt{Rejected}, \texttt{Promoted}) and seven authorization types (four reviewer verdicts, agent move, Zar approve, Zar reject). Key proved properties include:
\begin{itemize}[leftmargin=1.5em]
\item \texttt{Rejected} and \texttt{Promoted} are terminal---no valid outgoing transitions exist from either state,
\item promotion from staging requires \texttt{ReviewerAllow}; agent moves can never directly reach \texttt{Promoted} or \texttt{Rejected},
\item \texttt{PendingZar} can only exit via \texttt{ZarApprove} or \texttt{ZarReject}---the reviewer cannot bypass human escalation,
\item staging states cannot skip to \texttt{Promoted} via \texttt{ZarApprove}---the path must go through \texttt{PendingZar} first.
\end{itemize}

\subsection{Runtime Source Validation}
The \texttt{move\_existing\_artifact()} function in \texttt{security\_review.rs} now validates that the artifact's source path is under a valid mindlock subdirectory (\texttt{in/}, \texttt{out/}, \texttt{work/}, or \texttt{pending-zar/}) using a full path prefix check against the configured \texttt{mindlock\_dir}. This enforces the terminal-state invariant at the I/O boundary: artifacts in \texttt{rejected/} cannot be moved, and paths outside the mindlock root are rejected entirely.

\subsection{Cleanup Retention}
Rejected artifacts are cleaned up after 30 days via a systemd timer running daily at hour 23. The cleanup eligibility predicate is itself Verus-verified: only terminal artifacts older than the retention period are eligible, and staging or pending artifacts are proved ineligible regardless of age.

\section{Reviewer and Human Approval}
The architecture includes a reviewer LLM and a human escalation path.

For staged artifacts, the reviewer can return:
\begin{itemize}[leftmargin=1.5em]
\item \texttt{Allow}
\item \texttt{RequestRevision}
\item \texttt{RequireZarApproval}
\item \texttt{Reject}
\end{itemize}

The review layer is intentionally treated as an external oracle in the formal model. VeriCore can prove that if the reviewer returns a verdict, the corresponding state machine transition is enforced correctly. VeriCore does \emph{not} attempt to prove that the reviewer made the semantically best or morally correct decision.

Human approval remains a first-class path. In particular, \texttt{self\_escalate} is meant to preserve agency: Oruzi can ask for help directly without the reviewer first granting permission to ask.

\section{LLM Completions Bridge}
An earlier version of VeriCore made direct HTTP calls to OpenRouter from Rust, bypassing OpenClaw's model resolution chain. This created silent cross-provider billing, duplicated auth logic, and meant that session model overrides (the \texttt{/model} command) had no effect on VeriCore's driver turns.

The current architecture replaces that path entirely with a Unix socket bridge:
\begin{enumerate}[label=\arabic*.]
\item Rust's \texttt{GatewayLlmClient} sends an OpenAI-format request (4-byte big-endian length prefix + JSON) to a Unix socket.
\item The TypeScript completions bridge (\texttt{completions-server.ts}) receives the request and converts OpenAI message/tool formats to pi-ai types.
\item The bridge runs the request through OpenClaw's full model resolution chain: \texttt{resolveConfiguredModelRef} $\to$ \texttt{resolveStoredModelOverride} $\to$ \texttt{runWithModelFallback} $\to$ \texttt{resolveModel} $\to$ \texttt{getApiKeyForModel} $\to$ \texttt{complete()}.
\item The response is converted back to OpenAI format and returned over the socket.
\end{enumerate}

The key architectural invariant is that \textbf{Rust has zero model, provider, or auth configuration}. The legacy \texttt{model}, \texttt{api\_base}, and \texttt{api\_key\_env} fields have been purged from the Rust config. All provider routing, subscription-based auth profiles, fallback logic, and billing semantics live in OpenClaw's TypeScript runtime.

Both the driver turn LLM call and the security reviewer LLM call use the same bridge. The bridge distinguishes them via \texttt{call\_kind} (for logging) and \texttt{prompt\_mode} (for context), but both resolve through the same model/auth stack.

\subsection{Fallback Notifications}
When the bridge falls back to a different model or provider, the operator is notified via Telegram DM. This is not optional observability---it is an essential feature, because silent provider changes can alter billing, capability, and safety characteristics. The notification path uses the same \texttt{routeReply} mechanism as normal Telegram messages, with a 5-minute deduplication window per fallback transition.

The TypeScript side is the authoritative model-resolution ledger, writing structured JSONL audit entries for every completion.

\subsection{Lessons Learned}
The bridge implementation exposed several verification gaps worth documenting:
\begin{itemize}[leftmargin=1.5em]
\item The initial bridge was deployed with only infrastructure verification (socket exists, server starts) but no end-to-end LLM path test. A local test harness (\texttt{scripts/testbus.sh}) was needed to exercise the full VeriCore $\to$ bridge $\to$ provider path without requiring real Telegram DMs.
\item Legacy Rust config fields (\texttt{api\_base}, \texttt{api\_key\_env}, \texttt{model}) caused a silent routing bug: VeriCore appeared to work but was hitting OpenRouter directly (returning HTTP 402) instead of going through the bridge. The fix was to purge these fields entirely.
\item A string comparison for fallback detection (\texttt{configured\_model != resolved\_model}) produced false positives on every clean run because the two fields use different naming conventions. The fix was to gate on the boolean \texttt{fallback\_used} field instead.
\item The security reviewer path through the bridge has not yet been E2E tested---it is a known gap.
\end{itemize}

\section{What Is Verified Today}
The crate \texttt{vericore-policy} contains verified policy logic at two levels. As of March 2026, Verus verifies 116 functions with 0 errors across 10 modules.

\subsection{Verified Runtime Kernel}
The following modules are production-wired: the same functions that Verus checks are the ones \texttt{vericore-core} calls for live policy decisions.
\begin{itemize}[leftmargin=1.5em]
\item tier definitions and flow-label algebra (\texttt{tiers.rs}): context rank, integrity rank, flow-to, label join. Proved properties include join idempotence, commutativity, associativity, and the duality between \texttt{can\_flow\_to} and \texttt{can\_read}.
\item default channel-to-context and context-to-integrity mapping (\texttt{channels.rs}): called by \texttt{GatePolicy::context\_for\_channel} and \texttt{integrity\_for\_channel} in production. Proved: public channels map to \texttt{Untrusted}, private channels map to \texttt{Trusted}, no channel maps above \texttt{Private}, and the integrity mapping is injective.
\item ingress decision predicates (\texttt{ingress.rs}): \texttt{ingress\_allows\_exec} and \texttt{ingress\_allows\_write}, called from \texttt{check\_ingress\_integrity} in production. Proved: \texttt{Reviewed} cannot exec, \texttt{Trusted} has full capability, write permissions are monotone in integrity.
\end{itemize}

\subsection{Verified Reference Specs}
Seven additional modules provide proved reference specs. These are not runtime-called kernels; correspondence with runtime behavior is established by tests, not by the prover. Their value is as regression guards and precise invariant statements.

\paragraph{Model resolution} (\texttt{model\_resolution.rs}). Proves candidate-list policy properties over an abstract \texttt{Seq<ModelRef>}: primary-first ordering, no duplicates, cross-provider gating, and provider stability. Also provides executable decision functions for resolution metadata.

\paragraph{LLM bridge contract} (\texttt{llm\_bridge.rs}). Specifies the completions bridge interface: request well-formedness, response usability, call-kind/prompt-mode consistency (security review must use reviewer mode), and session-key requirements (driver turns require a session key).

\paragraph{Mindlock state machine} (\texttt{mindlock.rs}). The full transition relation, staging/terminal classification, source validation, and cleanup eligibility (see Section~5).

\paragraph{Receipt authorization} (\texttt{receipt.rs}). Promotion requires either a valid receipt (hash match AND target match) or explicit Zar approval with audit trail. Proved: tamper detection, target-swap detection, and the at-least-one-authorization invariant.

\paragraph{Egress flow control} (\texttt{egress.rs}). Information flows upward in secrecy: \texttt{Public} can flow to anything, \texttt{Private} can only flow to \texttt{Private}. Proved: transitivity, same-tier always allowed, private-to-public denied, and security review triggers.

\paragraph{Gate chain} (\texttt{gate\_chain.rs}). The runtime gate chain is a strict conjunction of six boolean gates (schedule, action kind, tool ID, skill ID, ingress integrity, primitive action). Proved: any single deny breaks the chain, and the chain is monotone.

\paragraph{Cross-module compositions} (\texttt{compositions.rs}). Six proofs chaining \texttt{channels} $\to$ \texttt{tiers} $\to$ \texttt{ingress}: public Telegram/API/Moltbook channels cannot execute commands, untrusted channels cannot write to private targets, and private channels have full capability.

\subsection{The Distinction}
This separation matters and should be maintained honestly:
\begin{itemize}[leftmargin=1.5em]
\item \textbf{Verified runtime kernel}: Verus-checked functions that production calls directly. A bug in these functions would be caught by the prover before deployment.
\item \textbf{Verified reference spec}: Verus-checked invariants that define what the runtime \emph{should} satisfy. Correspondence between spec and runtime is established by tests, not by the prover.
\item \textbf{Tested I/O and runtime behavior}: integration tests covering transport, filesystem, async races, and service availability. Verus does not and should not cover these.
\end{itemize}

\textbf{Remaining drift surface:} \texttt{vericore-core} still has its own copies of the \texttt{Channel}, \texttt{ContextTier}, and \texttt{IntegrityTier} enums, with total-match conversion helpers bridging them to the verified crate. This duplication is acceptable as bootstrap scaffolding but should eventually be consolidated so that verified types are used directly.

\section{Proof Boundary}
The current proof boundary covers all deterministic policy decision surfaces:
\begin{enumerate}[label=\arabic*.]
\item label algebra over secrecy and integrity (join lattice properties, flow/read duality),
\item channel-to-context and context-to-integrity mapping (injectivity, ceiling bounds),
\item ingress allow/deny decisions over normalized target classes (monotonicity, full capability),
\item mindlock state machine (terminal states, valid transitions, source validation, cleanup eligibility),
\item receipt and human-approval conditions for promotion (tamper detection, at-least-one-auth),
\item egress flow control (transitivity, private-to-public denial, review triggers),
\item gate chain conjunction (monotonicity, any-deny-breaks-chain),
\item bridge contract (call-kind/prompt-mode consistency, session-key requirements),
\item cross-module composition proofs (channel $\to$ ingress $\to$ capability).
\end{enumerate}

The following are explicitly \emph{outside} the proof boundary:
\begin{itemize}[leftmargin=1.5em]
\item Telegram UX,
\item async daemon behavior,
\item filesystem race handling in detail,
\item reviewer prompt quality,
\item RAG relevance,
\item social or stylistic policy guidance.
\end{itemize}

This is not a weakness. It is a scoping decision. Formal verification is being used where deterministic invariants matter most and where the code can realistically remain simple enough to verify.

\section{Current Weak Point: I/O Source Authenticity}
The architecture is strong against \emph{content spoofing} but weaker against \emph{source spoofing}.

Content spoofing means message text tries to impersonate a more trusted source. VeriCore is already robust against this because authority comes from the channel assigned by the transport, not from the text itself.

Source spoofing is different. If an attacker can submit a forged daemon request with a falsely labeled channel such as \texttt{telegram\_dm}, VeriCore will trust that label and derive \texttt{Private/Trusted}. That means the real trust boundary is not only the policy function; it is also the authenticity of the gateway-to-daemon envelope.

The right next hardening step is therefore transport authenticity:
\begin{itemize}[leftmargin=1.5em]
\item authenticate daemon requests,
\item restrict who can submit control/private stimuli,
\item later, sign or MAC file-based or agent-bridge envelopes.
\end{itemize}

\section{Why a Verified Kernel Crate}
Real verification-heavy systems do not stop at disconnected models forever. They either verify the actual code or a very thin executable kernel that the runtime actually calls. That is the design goal here.

The crate split is therefore pragmatic:
\begin{itemize}[leftmargin=1.5em]
\item keep async I/O, sockets, Telegram, and filesystem mutation in normal Rust; delegate LLM interaction to TypeScript via the completions bridge,
\item move pure deterministic decision functions into a crate that is both executable and Verus-verifiable,
\item have production call those functions directly.
\end{itemize}

This preserves runtime practicality while shrinking the gap between ``proved model'' and ``running system''.

\section{Next Steps}
The correct near-term sequence is:
\begin{enumerate}[label=\arabic*.]
\item \textbf{(done)} ingress predicates Verus-verified and production-wired,
\item \textbf{(done)} default channel/context/integrity mapping verified and production-wired,
\item \textbf{(done, reference spec)} model resolution candidate-list policy verified,
\item \textbf{(done)} completions bridge with fallback notifications and local E2E test harness,
\item \textbf{(done, reference spec)} mindlock state machine, receipt authorization, egress flow control, gate chain, bridge contract, and cross-module compositions verified,
\item \textbf{(done)} mindlock runtime hardening: source validation guard with full path prefix check; 30-day rejected cleanup via systemd timer,
\item E2E test the security reviewer path through the completions bridge (known gap),
\item wire \texttt{model\_resolution} into Rust runtime (currently reference-only; TS is authoritative),
\item wire \texttt{llm\_bridge} consistency checks as debug assertions,
\item harden the gateway-to-daemon authenticity boundary,
\item consolidate enum types so production uses verified types directly.
\end{enumerate}

Testing is still required. Verus does not replace integration tests for socket boundaries, filesystem behavior, async races, or service availability. The practical split is:
\begin{itemize}[leftmargin=1.5em]
\item \textbf{Verus (runtime kernel)}: pure policy functions called by production,
\item \textbf{Verus (reference spec)}: proved invariants mirrored by runtime tests,
\item \textbf{tests}: I/O, operational integration, and spec--runtime correspondence.
\end{itemize}

\section{Conclusion}
VeriCore is now past the point where detached proof sketches alone are the main value. 116 functions are verified with 0 errors across 10 modules spanning three runtime-wired kernel modules, seven reference spec modules, and cross-module composition proofs.

For the runtime-called kernel (tiers, channels, ingress), the verified predicates are the authoritative runtime logic. For the reference specs (model resolution, mindlock, receipt, egress, gate chain, bridge contract), the invariants are stated precisely and proved consistent, serving as regression guards and precise targets for future runtime wiring.

The mindlock state machine and receipt authorization---previously identified as next targets---are now fully specified and verified. Runtime enforcement has been strengthened: \texttt{move\_existing\_artifact()} validates source paths against the configured mindlock root, and rejected artifacts are cleaned up on a 30-day retention schedule.

That combination---verified deterministic policy for the runtime kernel, verified reference specs for adjacent surfaces, runtime enforcement of verified invariants, and tested transport and operational boundaries---is the strongest realistic path for this architecture.

\end{document}
